\documentclass[12pt]{article}
\usepackage[utf8]{inputenc}
\usepackage[T1]{fontenc}
\usepackage[margin=1in]{geometry}
\usepackage{setspace}
\usepackage{xcolor}
\usepackage[colorlinks=true, urlcolor=blue!60!black, linkcolor=blue!60!black]{hyperref}
\setlength{\parindent}{2em}
\setlength{\parskip}{1.25em}
\renewcommand{\baselinestretch}{1.0}
\begin{document}
\begin{center}
{\LARGE \textit{Chisholm v. Georgia}, 2 U.S. (2 Dall.) 419 (1793)\par}
\vspace{0.8em}
{\large Case Brief --- Four-Page Edition\par}
\end{center}
\vspace{1em}

\section{Memory Jogger}

In the Court's first major constitutional decision, a South Carolina executor sued Georgia on a Revolutionary War debt; the Justices held 4--1, in seriatim opinions, that Article III lets a citizen of one State sue another State without its consent---a ruling the Eleventh Amendment promptly overturned (\href{https://supreme.justia.com/cases/federal/us/2/419/}{2 U.S. 419}).

\section{Detailed Case Facts}

In 1777, Georgia's commissioners Thomas Stone and Edward Davies bought Revolutionary War supplies for American troops near Savannah from Robert Farquhar, a South Carolina merchant. Although Georgia's treasury gave the commissioners funds, they never paid Farquhar, who died in 1784 with the debt unpaid. Alexander Chisholm, a South Carolina citizen and executor of Farquhar's estate, petitioned the Georgia legislature for payment, which denied his claim in 1789.

\section{Procedural History}

Chisholm first sued in the federal circuit court in Georgia, where Governor Edward Telfair and the State's Attorney General refused to submit to jurisdiction. In 1791 Justice Iredell, riding circuit, and District Judge Nathaniel Pendleton dismissed for lack of jurisdiction. In 1792 Chisholm brought an action of \emph{assumpsit} directly in the Supreme Court under Article III and Section 13 of the Judiciary Act of 1789. Georgia refused to appear, and Attorney General Edmund Randolph, arguing alone in February 1793, moved for a conditional default judgment.

\section{Judicial Votes}

For the plaintiff, \textbf{4--1}, decided February 19, 1793.
\begin{itemize}
\item \textbf{Majority (seriatim):} Chief Justice Jay and Justices Cushing, Wilson, and Blair, each writing separately; there was no opinion of the Court.
\item \textbf{Dissent:} Justice Iredell.
\end{itemize}

\section{Holding}

Article III extends federal judicial power to controversies between a State and citizens of another State whether the State is plaintiff or defendant, and gives the Supreme Court original jurisdiction. Because sovereignty resides in the people of the United States, common-law sovereign immunity does not bar an out-of-state citizen's \emph{assumpsit} action against Georgia. Service on the Governor and Attorney General was valid, and default judgment would be entered unless Georgia appeared or showed cause by the next Term.

\section{Analysis of Opinions}

\textbf{Text of Article III (the majority).} Rejecting the ratification-era assurances of Hamilton in \emph{Federalist No.~81}, Madison, and Marshall that States could appear only as plaintiffs, Blair and Cushing reasoned that adjacent clauses necessarily make States defendants. Jay noted that ``party'' includes defendants and that Congress distinguished plaintiffs from defendants for ambassadors but not for States.

\textbf{Popular sovereignty (Wilson and Jay).} Wilson asked whether ``the people of the United States form a Nation''; royal immunity stemmed from feudal hierarchy, but here ``there are citizens, but no subjects,'' and ``As to the purposes of the Union, therefore, Georgia is NOT a sovereign state.'' Jay traced sovereignty from the Crown to the people and noted that Georgia was itself suing as a plaintiff in \emph{Georgia v. Brailsford}: a rule is bad that does not work both ways.

\textbf{Structure.} Wilson argued that the Contracts Clause would be hollow if States could breach contracts with impunity; Jay and Cushing stressed a neutral national tribunal for interstate disputes. Blair relied strictly on text, and Cushing observed that if the Constitution proved inconvenient, it provided a mode of amendment.

\textbf{Iredell (dissenting).} Federal courts cannot fashion remedies without Congress. Section 14 of the Judiciary Act limited writs to those ``agreeable to the principles and usages of law,'' and at common law no compulsory \emph{assumpsit} lay against the Crown; relief came by petition as ``a matter of grace, though not upon compulsion.'' Each State remains ``completely sovereign'' as to reserved powers, and he intimated opposition to ``a compulsive suit against a State for the recovery of money.''

\textbf{Aftermath.} Congress proposed the Eleventh Amendment in 1794 and it was ratified in 1795; in \emph{Hollingsworth v. Virginia} (1798) the Court dismissed pending suits against States. \emph{Hans}, \emph{Alden}, and \emph{Hyatt} extended immunity, while \emph{Ex parte Young} and \emph{Fitzpatrick} created exceptions.

\section{Examples: Future Applications}

\begin{enumerate}
\item \textbf{Same-Side: The State as Plaintiff.} Louisiana sues a Texas supplier for breach of warranty. The Eleventh Amendment reaches only suits ``commenced or prosecuted against'' a State, the asymmetry Jay criticized.
\item \textbf{Same-Side: Prospective Relief.} A supplier sues the state treasurer to enjoin a statute cancelling state contracts as a Contracts Clause violation. \emph{Ex parte Young} allows it, though not payment of the past debt.
\item \textbf{Opposite-Side: Unpaid Contract Debt.} A Texas supplier sues Louisiana in federal court for \$15 million. \emph{Chisholm} would allow it; the Eleventh Amendment and \emph{Hans} now bar it.
\item \textbf{Opposite-Side: State-Court Wage Claims.} Maine probation officers sue the State in its own courts for federal overtime. \emph{Alden v.\ Maine} bars the suit on structural grounds traced to Iredell.
\item \textbf{Fence-Sitter: Hospital Authority.} A separately incorporated hospital authority is responsible for its own debts, but the treasury must pay judgments above a threshold. Under \emph{Galette} (2026), separate incorporation weighs against immunity while state liability weighs for it.
\end{enumerate}

\section{Critique}

\textbf{Textualism inverted.} The \emph{Chisholm} majority read Article III literally, while modern sovereign immunity rests on unwritten structural postulates drawn from Iredell's dissent, which itself rested formally on the Judiciary Act of 1789.

\textbf{Diversity versus immunity theories.} Critics argue that the Eleventh Amendment merely repealed Citizen--State diversity jurisdiction, leaving federal-question suits against States intact.

\section{Key Quotations}

\begin{quote}
``No such ideas obtain here; at the Revolution, the sovereignty devolved on the people, and they are truly the sovereigns of the country, but they are sovereigns without subjects (unless the African slaves among us may be so called), and have none to govern but themselves; the citizens of America are equal as fellow citizens, and as joint tenants in the sovereignty.''
\end{quote}
\noindent --- Chief Justice Jay, 2 U.S. at 471--472.

\begin{quote}
``Every State in the Union, in every instance where its sovereignty has not been delegated to the United States, I consider to be as completely sovereign as the United States are in respect to the powers surrendered. The United States are sovereign as to all the powers of Government actually surrendered: each State in the Union is sovereign as to all the powers reserved.''
\end{quote}
\noindent --- Justice Iredell, dissenting, 2 U.S. at 435.

\begin{quote}
``A state, like a merchant, makes a contract. A dishonest state, like a dishonest merchant, wilfully refuses to discharge it. The latter is amenable to a court of justice. Upon general principles of right, shall the former, when summoned to answer the fair demands of its creditor, be permitted, Proteus-like, to assume a new appearance, and to insult him and justice by declaring `I am a Sovereign state?' Surely not.''
\end{quote}
\noindent --- Justice Wilson, 2 U.S. at 456.

\end{document}
